\chapter{Baseline Reproduction and Adaptation}
\label{chap:replication}

\todo[inline]{Outline approved; prose pending. Present the baseline verification and adaptation as evidence-based experimental work. Report unresolved discrepancies explicitly and keep the original and dense evaluation protocols separate.}

\section{Initial reproduction attempts with TranAD}
\label{sec:replication-tranad}
% Document concrete attempts, available configurations and observed difficulties.
% Avoid unsupported claims that a published method is intrinsically irreproducible.

\section{TranAD+ implementation and configuration reconstruction}
\label{sec:replication-tranadplus}
% Compare local/reference sources and recover configurations from official run logs.
% Identify which run supplies each published POT/BF result, including different winners.

\section{Reproduction of the original benchmarks}
\label{sec:replication-benchmarks}
% Compare the six datasets with matched configurations and original evaluation settings.
% Preserve differing batch sizes/configurations as separate experiments.
% Put full per-run configurations and supplementary tables in Appendix A.

\section{Sensitivity to the threshold-search protocol}
\label{sec:replication-thresholds}
% Compare original and dense searches on the same checkpoint and scores.
% Attribute changes caused by evaluation to the protocol rather than model improvements.

\section{Adaptation of TranAD+ to OPS-SAT}
\label{sec:replication-opssat}
% Describe changes required for single-feature fragments and the observed baseline results.
% Keep fragment validation/test outcomes distinct from point-adjusted oracle metrics.

\section{Reproduction assessment and remaining discrepancies}
\label{sec:replication-assessment}
% Summarize verified agreement, explained deviations and unresolved differences.
% Establish the reference configurations used by the following ablation study.
